\section{Related Work}
\paragraph{Feedback-conditioned memory.} Prior work covers reward reliability~\citep{wang2026planrewardbench}, success/failure-derived memory~\citep{ouyang2026reasoningbank}, reward-conditioned construction~\citep{wang2025memalpha}, and scoring/reuse distortion~\citep{asadolahi2026inflation,yang2026romerl}. These are not our novelty; the frozen writer branch is only our controlled intervention.

\paragraph{Retrieval, utilization, and our residual.} Mem2ActBench studies active use~\citep{shen2026mem2actbench}; QCR separates retrieval from reuse~\citep{li2026qcr}; Yuan et al. diagnose write/retrieval/utilization~\citep{yuan2026retrievalutilization}; and lifecycle work names the same broad stages~\citep{lin2026memorylifecycle}. These concepts are not our novelty. Our residual question is intervention-specific: after the same source trajectory is converted into two branch-conditioned persistent states, which measured boundary still carries that branch contrast? The load-bearing distinction is separation of exposure from uptake while preserving the identity of the controlled branch intervention, not lifecycle naming or a power-matched causal bottleneck. A retrieval hit can establish source-item availability without establishing transport of the branch-differentiating residual into the measured decision. Known self-improvement variance provides complementary motivation~\citep{ye2026fragility}.
